\documentclass[11pt]{article}
\usepackage[utf8]{inputenc}
\usepackage[T1]{fontenc}
\usepackage{geometry}
\usepackage{enumitem}
\usepackage{amsmath}
\usepackage{amssymb}
\usepackage{xcolor}
\usepackage{hyperref}
\usepackage{booktabs}
\usepackage{array}

\geometry{margin=1in}

\title{\textbf{Midterm Assignment: Prompt Engineering}\\
\large Weeks 1--5 Comprehensive Assessment}
\author{Prompt Engineering Course\\Ankara Medipol University}
\date{Due Date: 15.11.2025 until 23.59}

\begin{document}

\maketitle

\section*{Assignment Overview}

\textbf{Total Points:} 100\\
\textbf{Submission Format:} Single PDF document\\
\textbf{Due Date:} 15.11.2025 until 23.59\\
\textbf{Learning Objectives:} Demonstrate mastery of fundamental prompt engineering concepts, techniques, and frameworks covered in Weeks 1--5.

\section*{Important Instructions}

\begin{itemize}
    \item Submit \textbf{ONE PDF file} containing all parts of the assignment.
    \item Use any AI tool available to you (ChatGPT, Claude, Gemini, etc.) --- specify which tool you used.
    \item Include \textbf{screenshots} of your prompts and AI responses for verification.
    \item Label each section clearly with part numbers (Part 1, Part 2, etc.).
    \item Follow word count limits strictly (±10\% acceptable).
    \item \textbf{Academic integrity:} Your prompts must be your own work. Explain your reasoning clearly.
\end{itemize}

\section*{Assignment Structure}

This assignment consists of 3 parts:

\begin{enumerate}
    \item \textbf{Prompt Comparison Analysis} (30 points)
    \item \textbf{Practical Prompt Design} (40 points)
    \item \textbf{Critical Reflection} (30 points)
\end{enumerate}

\vspace{0.5cm}

\hrule

\section{Part 1: Prompt Comparison Analysis (30 points)}

\subsection*{Task}

Select ONE domain from the following: \textbf{Healthcare, Legal, Marketing, Education, or Engineering}.

For your chosen domain, create TWO versions of a prompt to accomplish the same task:

\begin{itemize}
    \item \textbf{Version A:} A basic, vague prompt
    \item \textbf{Version B:} An improved prompt using advanced prompting techniques
\end{itemize}

Then, test BOTH prompts with an AI tool and analyze the results.

\subsection*{Deliverables}

\begin{enumerate}[label=\alph*.]
    \item \textbf{Domain \& Task Description} (50 words max)
    \begin{itemize}
        \item State your chosen domain
        \item Describe the specific task (e.g., ``Draft a patient education handout about diabetes'')
    \end{itemize}
    
    \item \textbf{Version A: Vague Prompt}
    \begin{itemize}
        \item Write the vague prompt
        \item Include screenshot of AI response
    \end{itemize}
    
    \item \textbf{Version B: Improved Prompt}
    \begin{itemize}
        \item Write the improved prompt
        \item \textbf{Label which techniques you used} (e.g., role prompting, delimiters, few-shot examples, chain-of-thought, structured output)
        \item Include screenshot of AI response
    \end{itemize}
    
    \item \textbf{Comparison Analysis} (200--250 words)
    \begin{itemize}
        \item Compare the quality, usefulness, and specificity of both outputs
        \item Identify at least 3 concrete improvements in Version B's output
        \item Explain WHY the improved prompt worked better (reference course concepts)
    \end{itemize}
\end{enumerate}

\hrule
\vspace{0.5cm}

\section{Part 2: Practical Prompt Design (40 points)}

\subsection*{Task}

Design a complete prompt using \textbf{ONE of the following frameworks}:

\begin{itemize}
    \item CO-STAR Framework (Context, Objective, Style, Tone, Audience, Response)
    \item RACI Prompt Canvas (Role, Audience, Context, Instructions)
    \item IDEA Loop (Instruct, Detail, Example, Adjust)
    \item SMART Prompt Outputs (Specific, Measurable, Actionable, Relevant, Time-bound)
\end{itemize}

\textbf{Scenario:} You are helping a small business owner create content. Choose ONE of these tasks:

\begin{enumerate}
    \item Design a customer onboarding email sequence (3 emails)
    \item Create a social media content calendar for one month
    \item Draft a competitive analysis report for a new market
    \item Generate a product launch announcement strategy
\end{enumerate}

\subsection*{Deliverables}

\begin{enumerate}[label=\alph*.]
    \item \textbf{Framework Selection \& Justification} (100 words)
    \begin{itemize}
        \item State which framework you chose and which business task
        \item Explain WHY this framework is appropriate for this task
    \end{itemize}
    
    \item \textbf{Complete Prompt Following Framework Structure}
    \begin{itemize}
        \item Write your full prompt
        \item \textbf{Clearly label each framework component} (e.g., ``[CONTEXT]:'', ``[OBJECTIVE]:'', etc.)
        \item Ensure all framework elements are present
    \end{itemize}
    
    \item \textbf{AI Output}
    \begin{itemize}
        \item Include screenshot or full text of AI response
        \item Highlight or annotate where the AI successfully followed your framework
    \end{itemize}
    
    \item \textbf{Refinement \& Iteration} (150--200 words)
    \begin{itemize}
        \item Identify at least 2 weaknesses or areas for improvement in the initial output
        \item Write ONE refined follow-up prompt that addresses these issues
        \item Include screenshot of the improved output
        \item Explain what changed and why
    \end{itemize}
    
    \item \textbf{Framework Effectiveness Reflection} (150 words)
    \begin{itemize}
        \item Evaluate how well the framework helped you structure the prompt
        \item Would you use a different framework for this task? Why or why not?
        \item What did you learn about framework selection?
    \end{itemize}
\end{enumerate}

\hrule
\vspace{0.5cm}

\section{Part 3: Critical Reflection (30 points)}

\subsection*{Task}

Synthesize your learning from the first half of the course by reflecting on your prompt engineering journey so far.

\subsection*{Deliverables}

Write a structured reflection (400--500 words total) addressing the following:

\begin{enumerate}[label=\alph*.]
    \item \textbf{Skill Progression} (150 words)
    \begin{itemize}
        \item How has your understanding of prompt engineering evolved in the first half of this course?
        \item Identify one specific ``aha moment'' or breakthrough in your learning so far
        \item What was the most challenging concept and how did you overcome it?
    \end{itemize}
    
    \item \textbf{Practical Application} (150 words)
    \begin{itemize}
        \item Describe ONE real-world scenario where you could apply prompt engineering in your field/career
        \item Which frameworks or techniques would you use and why?
        \item What challenges do you anticipate and how would you address them?
    \end{itemize}
    
    \item \textbf{Limitations \& Ethics} (100--150 words)
    \begin{itemize}
        \item Identify 2 limitations of AI/LLMs that you must consider when prompting
        \item Discuss one ethical consideration in your chosen application scenario
        \item How would you ensure responsible use of AI in your prompts?
    \end{itemize}
    
    \item \textbf{Goals for the Second Half} (50 words)
    \begin{itemize}
        \item What aspects of prompt engineering do you want to explore further in the remaining course?
        \item What skill gap do you want to close by the end of the course?
    \end{itemize}
\end{enumerate}

\hrule
\vspace{0.5cm}

\section*{Submission Checklist}

Before submitting, ensure your PDF includes:

\begin{itemize}
    \item[$\square$] Cover page with your name, student ID, and submission date
    \item[$\square$] Part 1: Domain, both prompts, screenshots, analysis
    \item[$\square$] Part 2: Framework justification, labeled prompt, outputs, refinement, reflection
    \item[$\square$] Part 3: Critical reflection (400--500 words)
    \item[$\square$] Clear labeling of all sections and parts
    \item[$\square$] All screenshots are readable and properly captioned
    \item[$\square$] AI tool(s) used are specified
    \item[$\square$] Word counts are within acceptable range (±10\%)
    \item[$\square$] Document is formatted professionally (consistent fonts, spacing, headings)
\end{itemize}

\section*{Grading Summary}

\begin{center}
\begin{tabular}{llr}
\toprule
\textbf{Part} & \textbf{Description} & \textbf{Points} \\
\midrule
Part 1 & Prompt Comparison Analysis & 30 \\
Part 2 & Practical Prompt Design & 40 \\
Part 3 & Critical Reflection & 30 \\
\midrule
\textbf{Total} & & \textbf{100} \\
\bottomrule
\end{tabular}
\end{center}

\section*{Tips for Success}

\begin{itemize}
    \item \textbf{Start early:} Don't wait until the last day to test your prompts.
    \item \textbf{Iterate:} Try multiple versions of your prompts and choose the best ones.
    \item \textbf{Be specific:} Use exact terminology from the course when explaining your choices.
    \item \textbf{Show your work:} Screenshots and detailed explanations demonstrate your process.
    \item \textbf{Proofread:} Check for clarity, completeness, and professional presentation.
    \item \textbf{Reference course concepts:} Explicitly mention techniques, frameworks, and principles covered in the course.
\end{itemize}

\vspace{1cm}

\vspace{0.5cm}

\noindent\textbf{Good luck! This assignment is your opportunity to demonstrate everything you've learned about prompt engineering.}

\end{document}
